% from ega0-19-fr.tex
\oldpage[0\textsubscript{IV-1}]{96}

  comme $C$ est une $A$-algèbre formellement lisse, cet homomorphisme se
  factorise en
  $C\xrightarrow{f_\alpha}E_{\alpha,n}\to
  E_{\alpha,n}/(\mathfrak J_\alpha/\mathfrak J_\alpha^{n+1})$, où $f_\alpha$
  est un $A$-homomorphisme continu~; par suite,
  $\mathfrak J_\alpha/\mathfrak J_\alpha^{n+1}$ devient, au moyen de
  $f_\alpha$, un $C$-module topologique discret annulé par un idéal ouvert de
  $C$. L'hypothèse que $\mathfrak J/\mathfrak J^2$ est un $C$-module
  formellement projectif entraîne alors l'existence d'une application
  $C$-linéaire continue
  $g_\alpha:\mathfrak J/\mathfrak J^2\to
  \mathfrak J_\alpha/\mathfrak J_\alpha^{n+1}$ rendant commutatif le diagramme
  \[
    \xymatrix{
      \mathfrak J_\alpha/\mathfrak J_\alpha^{n+1} \ar[r] &
      \mathfrak J_\alpha/\mathfrak J_\alpha^2 \\
      & \mathfrak J/\mathfrak J^2 \ar[ul]^{g_\alpha} \ar[u]
    }
  \]
  Comme $f_\alpha$ et $g_\alpha$ sont continus, il y a un
  $\beta\geq\alpha$ tel que ces homomorphismes se factorisent respectivement en
  \[
    C\to C_\beta\xrightarrow{f'_{\alpha\beta}}E_{\alpha,n},
  \]
  \[
    \mathfrak J/\mathfrak J^2\to
    \mathfrak J_\beta/\mathfrak J_\beta^2
    \xrightarrow{g'_{\alpha\beta}}
    \mathfrak J_\alpha/\mathfrak J_\alpha^{n+1}
  \]
  et à partir de $f'_{\alpha\beta}$ et $g'_{\alpha\beta}$, on obtient donc
  canoniquement un homomorphisme de $C_\beta$-algèbres
  \[
    \mathbf S_{C_\beta}^*(\mathfrak J_\beta/\mathfrak J_\beta^2)
    \longrightarrow E_{\alpha,n}
  \]
  qui d'ailleurs (par définition de $g'_{\alpha\beta}$) s'annule dans la
  puissance $(n+1)$-ième de l'idéal d'augmentation de
  $\mathbf S_{C_\beta}^*(\mathfrak J_\beta/\mathfrak J_\beta^2)$~; passant au
  quotient par cette puissance $(n+1)$-ième, on obtient l'homomorphisme
  $v_{\alpha\beta}$ cherché, compte tenu des définitions de $f_\alpha$ et
  $g_\alpha$~; la surjectivité de $v_{\alpha\beta}$ résulte en effet de celle
  des deux homomorphismes $\gr^0(v_{\alpha\beta})$ et
  $\gr^1(v_{\alpha\beta})$, puisque cela entraîne que
  $\gr(v_{\alpha\beta})$ est surjectif (l'algèbre
  $\gr^*(E_{\alpha,n})$ étant engendrée par $\gr^0(E_{\alpha,n})$ et
  $\gr^1(E_{\alpha,n})$), et comme les filtrations considérées sont finies, on
  peut appliquer Bourbaki, \emph{Alg. comm.}, chap.~III, §~2,
  n\textsuperscript{o}~8, cor.~3 du th.~1.

  \item[\rm(ii)] L'hypothèse que $B$ est \emph{préadmissible} signifie que
  l'on peut supposer tous les $\mathfrak b_\alpha$ contenus dans un même
  $\mathfrak b_{\alpha_0}$ dont les puissances \emph{tendent vers} $0$. Cela
  entraîne en particulier que le noyau de l'homomorphisme de transition
  $E_{\gamma,n}\to E_{\alpha,n}$, égal à
  $(\mathfrak b_\alpha+\mathfrak J^{n+1})/
  (\mathfrak b_\gamma+\mathfrak J^{n+1})$, est \emph{nilpotent}~; appliquant le
  lemme (19.3.10.1), on voit qu'on peut supposer $f_\gamma$ choisi de sorte que
  le diagramme
  \[
    \xymatrix{
      & E_{\alpha,n} \\
      C \ar[ur]^{f_\alpha} \ar[r]_{f_\gamma} & E_{\gamma,n} \ar[u]
    }
  \]
  (où la flèche verticale est l'homomorphisme de transition) soit commutatif.
  L'hypothèse que $\mathfrak J/\mathfrak J^2$ est un $C$-module formellement
  projectif permet d'autre part de choisir
  $g_\gamma:\mathfrak J/\mathfrak J^2\to
  \mathfrak J_\gamma/\mathfrak J_\gamma^{n+1}$ de sorte que le diagramme
  \[
    \xymatrix@C=5em@R=3em{
      \mathfrak J_\alpha/\mathfrak J_\alpha^{n+1} \ar[r] &
      \mathfrak J_\alpha/\mathfrak J_\alpha^2 \\
      \mathfrak J_\gamma/\mathfrak J_\gamma^{n+1} \ar[u] \ar[r] &
      \mathfrak J_\gamma/\mathfrak J_\gamma^2 \ar[u] \\
      & \mathfrak J/\mathfrak J^2
        \ar[ul]_{g_\gamma} \ar[u]_{p_\gamma} \ar[uul]^{g_\alpha}
    }
  \]
  soit commutatif (il suffit de remarquer que l'image par $(g_\alpha,p_\gamma)$
  de $\mathfrak J/\mathfrak J^2$ dans le module produit
  $(\mathfrak J_\alpha/\mathfrak J_\alpha^{n+1})
  \times(\mathfrak J_\gamma/\mathfrak J_\gamma^2)$ est contenue (vu la
  définition de $g_\alpha$ et de la relation $\beta\leq\gamma$) dans l'image
  canonique $Q$ dans ce module produit du module
  $P=\mathfrak J_\gamma/\mathfrak J_\gamma^{n+1}$, et d'appliquer la définition
  (19.2.1) à l'homomorphisme \emph{surjectif} $P\to Q$ et à l'homomorphisme
  $(g_\alpha,p_\gamma)$ de $\mathfrak J/\mathfrak J^2$ dans $Q$). Ce choix de
  $f_\gamma$ et de $g_\gamma$ permet alors, en construisant
  $v_{\gamma\delta}$ comme dans \emph{(i)}, d'obtenir en outre la commutativité
  du diagramme (19.5.6.4).
\end{enumerate}
\end{env}

